Este anejo recoge los comandos mínimos necesarios para reconstruir la evidencia
experimental usada en la memoria. La intención no es sustituir al
\texttt{README.md}, sino dejar una ruta de ejecución compacta y verificable:
preparar el entorno, comprobar el paquete, generar datos, ejecutar campañas,
consolidar resultados y regenerar figuras.

\subsubsection*{Preparación del entorno}

El proyecto se gestiona con \texttt{uv}. La primera orden sincroniza el entorno
definido por \texttt{pyproject.toml} y \texttt{uv.lock}; la segunda ejecuta las
pruebas disponibles antes de lanzar campañas largas.

\begin{verbatim}
uv sync
uv run pytest
\end{verbatim}

\subsubsection*{Ejecución mínima de un escenario}

La ejecución elemental toma un escenario YAML, instancia vehículo, trayectoria,
controlador, perturbaciones y condiciones de terminación, y escribe dos
artefactos: \texttt{telemetry.json} con la serie temporal y
\texttt{metrics.json} con las métricas agregadas.

\begin{verbatim}
uv run simulador-quad run scenarios\hover_clean.yaml --no-visualization

uv run simulador-quad plot ^
  results\hover_clean\telemetry.json ^
  --metrics results\hover_clean\metrics.json ^
  --out results\hover_clean\figures_report ^
  --profile report --formats png pdf
\end{verbatim}

El mismo patrón permite ejecutar casos diagnósticos más exigentes:

\begin{verbatim}
uv run simulador-quad run scenarios\circle_noisy_wind.yaml --no-visualization
uv run simulador-quad run scenarios\lissajous_clean.yaml --no-visualization
uv run simulador-quad run scenarios\waypoint_clean.yaml --no-visualization
uv run simulador-quad run scenarios\composite_ood.yaml --no-visualization
\end{verbatim}

\subsubsection*{Dataset clásico y expertos PD}

El dataset clásico fija escenarios reproducibles y permite ajustar o congelar
controladores PD por familia. Tras modificar ganancias o criterios de ajuste,
los escenarios deben regenerarse con \texttt{--overwrite} para que las ganancias
embebidas y los manifiestos vuelvan a ser coherentes.

\begin{verbatim}
uv run python tools\generate_classic_dataset.py ^
  --version v1 --out data\classic_dataset\v1 --overwrite

uv run python tools\tune_classic_pid.py ^
  --dataset data\classic_dataset\v1 ^
  --out data\classic_dataset\v1\pids --workers 8

uv run python tools\run_classic_dataset.py ^
  --dataset data\classic_dataset\v1 ^
  --no-visualization --workers 4

uv run python tools\summarize_classic_dataset.py ^
  --dataset data\classic_dataset\v1
\end{verbatim}

\subsubsection*{Dataset neuronal de fuerza externa}

El controlador neuronal principal aprende por imitación la fuerza deseada del
lazo externo. Para ello se genera primero un banco de candidatos PD externos,
se selecciona el experto seguro por escenario y se exportan las muestras de
entrada y salida usadas por MLP, GRU y LSTM.

\begin{verbatim}
uv run python tools\generate_outer_force_pid_bank.py ^
  --dataset data\classic_dataset\v1 ^
  --out data\outer_force_pid_bank\v1 --workers 8

uv run python tools\generate_outer_force_dataset.py ^
  --source-dataset data\classic_dataset\v1 ^
  --pid-bank data\outer_force_pid_bank\v1 ^
  --out data\outer_force_dataset\v1

uv run python tools\train_neural_controller.py ^
  --dataset data\outer_force_dataset\v1 ^
  --architecture mlp ^
  --feature-version outer_force_min_v1 ^
  --out data\neural_control\outer_force_mlp_min_v1 ^
  --device auto

uv run python tools\evaluate_neural_controller.py ^
  --dataset data\outer_force_dataset\v1 ^
  --run data\neural_control\outer_force_mlp_min_v1 ^
  --device auto
\end{verbatim}

Las arquitecturas recurrentes se entrenan cambiando
\texttt{--architecture gru} o \texttt{--architecture lstm}. En la evaluación
cerrada se pasan siempre el checkpoint y la normalización correspondientes para
evitar mezclar modelos y escalados.

\subsubsection*{Batería OOD y comparación cerrada}

La batería OOD usada en resultados procede de
\texttt{data/neural\_ood/battery\_v1}. La auditoría de hélices fijó
\texttt{max\_duration\_s}=60\,s en el generador para que las trayectorias
trapezoidales rápidas dispongan de margen suficiente.

\begin{verbatim}
uv run python tools\generate_ood_battery.py ^
  --out data\neural_ood\battery_v1 ^
  --pid-source-dataset data\classic_dataset\v1 ^
  --overwrite

uv run python tools\run_classic_transfer_dataset.py ^
  --dataset data\neural_ood\battery_v1 ^
  --pid-source-dataset data\classic_dataset\v1 ^
  --pid-family all --include-native --workers 8 ^
  --no-visualization --rerun

uv run python tools\run_neural_outer_force_dataset.py ^
  --dataset data\neural_ood\battery_v1 --split ood ^
  --architecture mlp --device auto --workers 4 ^
  --no-visualization --rerun

uv run python tools\run_neural_outer_force_dataset.py ^
  --dataset data\neural_ood\battery_v1 --split ood ^
  --architecture gru --device auto --workers 4 ^
  --no-visualization --rerun

uv run python tools\run_neural_outer_force_dataset.py ^
  --dataset data\neural_ood\battery_v1 --split ood ^
  --architecture lstm --device auto --workers 4 ^
  --no-visualization --rerun

uv run python tools\summarize_comparison.py
\end{verbatim}

\subsubsection*{Figuras de resultados y anejos}

Las figuras comparativas del capítulo de resultados se regeneran desde el CSV
consolidado. Las figuras de telemetría de los anejos se generan desde
ejecuciones concretas para mostrar señales internas de simulación.

\begin{verbatim}
uv run simulador-quad plot-comparison ^
  results\comparison_all_runs.csv ^
  --out TFG_Memoria\Figuras\resultados ^
  --formats pdf png

uv run simulador-quad plot ^
  results\circle_noisy_wind\telemetry.json ^
  --metrics results\circle_noisy_wind\metrics.json ^
  --out TFG_Memoria\Figuras\resultados\sim_circle_noisy_wind ^
  --profile report --formats pdf png
\end{verbatim}

El mismo comando de \texttt{plot} se ha aplicado a \texttt{hover\_clean},
\texttt{lissajous\_clean}, \texttt{waypoint\_clean}, \texttt{circle\_drag} y
\texttt{composite\_ood}. En todos los casos, el perfil \texttt{report} elimina
títulos diagnósticos redundantes y conserva ejes, unidades, leyendas y formato
apto para memoria.

\subsubsection*{Artefactos de evidencia}

\begin{itemize}
  \item \texttt{results/comparison\_all\_runs.csv}: corridas comparables usadas
    por las figuras principales de resultados.
  \item \texttt{results/comparison\_summary.csv}: agregación por controlador,
    split y familia.
  \item \texttt{results/evidence\_manifest.csv}: procedencia de datasets,
    checkpoints y decisiones de inclusión.
  \item \texttt{results/*/telemetry.json}: series temporales por episodio.
  \item \texttt{results/*/metrics.json}: métricas y metadatos de cada
    ejecución.
  \item \texttt{TFG\_Memoria/Figuras/resultados/}: figuras finales en PDF y
    PNG.
\end{itemize}
